\documentclass[11pt,a4paper]{article}
\usepackage{brainteaser}

\begin{document}
\brainteasertitle{11}{Les 100 portes}

\begin{statementbox}
Cent portes sont initialement fermées. On effectue cent passages. Au passage
$1$, on change l'état de toutes les portes; au passage $2$, celui des portes
multiples de $2$; et, plus généralement, au passage $k$, l'état de chaque
porte dont le numéro est divisible par $k$.

Quelles portes sont ouvertes après le centième passage, et combien y en a-t-il?
\end{statementbox}

\begin{hintsbox}
\begin{enumerate}
  \item Fixer une porte numérotée $n$. À quels passages change-t-elle d'état?
  \item Quel lien existe-t-il avec le nombre de diviseurs de $n$?
  \item Pourquoi les diviseurs se regroupent-ils généralement par paires
  $d\leftrightarrow n/d$?
  \item Dans quel cas une paire se réduit-elle à un seul entier?
\end{enumerate}
\end{hintsbox}

\correction

La porte $n$ change d'état exactement lors des passages $k$ tels que $k\mid n$.
Elle est donc basculée $\tau(n)$ fois, où $\tau(n)$ est le nombre de diviseurs
positifs de $n$. Comme elle commence fermée, elle termine ouverte si et
seulement si $\tau(n)$ est impair.

Les diviseurs se regroupent par paires
\[
d\longleftrightarrow\frac nd.
\]
Tous sont appariés, sauf lorsqu'un diviseur vérifie
\[
d=\frac nd,
\]
c'est-à-dire $d^2=n$. Ainsi, $n$ possède un nombre impair de diviseurs si et
seulement si $n$ est un carré parfait.

\begin{resultbox}
Les portes ouvertes sont
\[
\boxed{1,4,9,16,25,36,49,64,81,100}.
\]
Il y en a $\boxed{10}$.
\end{resultbox}

\end{document}
